\begin{abstract}
全景语义分割在自动驾驶等领域具有重要应用价值，但现有基于视觉大模型的方法面临训练成本高、边界不连续等挑战。本文提出一种高效的全景分割框架，包含三项核心创新：(1) 记忆Token压缩(MTC)通过空间聚合将历史记忆规模压缩8倍；(2) 稀疏记忆检索(SMR)基于重要性动态选择关键记忆，进一步降低计算开销；(3) 无缝循环记忆(SCM)通过循环填充解决360°全景图边界断裂问题。在Cityscapes到DensePASS的无监督域适应任务中，我们的方法在保持精度的同时显著提升效率：8倍Token压缩率、55.8\%注意力时间降低，仅损失0.24\%的mIoU (60.22$\rightarrow$60.00，5-seed验证)。SCM在推理阶段提升0.33\%整体精度，交通标志类别提升5.04\% (26.07$\rightarrow$31.11)，验证了边界连续性建模的有效性。多随机种子统计验证表明方法具有良好稳定性，为资源受限场景下的全景分割应用提供了实用解决方案。
\end{abstract}

\noindent\textbf{关键词：} 全景语义分割；记忆优化；边界感知；无监督域适应；高效推理
